% Figure: the falling-sphere (Hoeppler-type) viscometer. A vertical tube
% of clear fluid with a sphere descending at its terminal velocity between
% two timing marks a distance L apart; the force balance on the sphere
% (weight down, buoyancy and Stokes drag up) is drawn alongside.
\begin{figure}[htbp]
\centering
\begin{tikzpicture}[>=stealth, scale=1.0]
  % tube
  \draw[engblack, very thick] (0,0) rectangle (2.2,8);
  \fill[engblue!12] (0.05,0.05) rectangle (2.15,7.95);
  % timing marks
  \draw[engred, thick] (0,6) -- (2.2,6) node[right, font=\scriptsize, engred] {upper mark};
  \draw[engred, thick] (0,1.6) -- (2.2,1.6) node[right, font=\scriptsize, engred] {lower mark};
  \draw[engred, <->] (-0.4,1.6) -- (-0.4,6) node[midway, left, font=\small] {$L$};
  % sphere descending
  \shade[ball color=enggray] (1.1,3.8) circle (0.5);
  \node[font=\scriptsize] at (1.1,3.8) {$a,\rho_s$};
  % velocity arrow
  \draw[engblack, very thick, ->] (1.1,3.2) -- (1.1,2.3)
    node[right, font=\small] {$U$};
  % force-balance free body to the right
  \begin{scope}[xshift=5.2cm, yshift=3.3cm]
    \shade[ball color=enggray] (0,0) circle (0.55);
    \draw[engred, very thick, ->] (0,-0.55) -- ++(0,-1.35)
      node[below, font=\scriptsize] {weight $\tfrac{4}{3}\pi a^3\rho_s g$};
    \draw[engblue, very thick, ->] (-0.25,0.55) -- ++(0,1.35)
      node[above left, font=\scriptsize] {buoyancy};
    \draw[englime, very thick, ->] (0.25,0.55) -- ++(0,1.0)
      node[above right, font=\scriptsize] {$6\pi\mu a U$};
  \end{scope}
\end{tikzpicture}
\caption[Falling-sphere viscometer]{The falling-sphere viscometer reads
the viscosity of a clear fluid from the terminal velocity of a sphere
descending between two timing marks a distance $L$ apart. At terminal
velocity the weight balances the sum of buoyancy and Stokes drag, so the
measured speed $U=L/t$ inverts directly to the viscosity through
$\mu=2a^2(\rho_s-\rho_f)g/(9U)$. The reading is valid only while the
sphere Reynolds number stays below unity, which sets an upper bound on
the sphere size and a lower bound on the fluid viscosity the instrument
can handle.}
\label{fig:vol03ch12:falling-sphere}
\end{figure}
